\section*{Bloque D. Ampliación: Teoría y Aplicaciones Prácticas (70\% / 30\%)}

\noindent\textit{Preguntas 91 a 180. Ampliación del banco oficial con 63 preguntas teóricas y conceptuales (91 a 153) y 27 aplicaciones cuantitativas prácticas (154 a 180).}

\subsection*{Parte I. Preguntas teóricas y conceptuales (91 a 153)}

\q{Si un monopolista maximizador de beneficios enfrenta una demanda con elasticidad precio constante igual a $\varepsilon = -0.8$ y un costo marginal estrictamente positivo, ¿qué puede afirmarse sobre su equilibrio?}
\begin{opts}
    \item No existe un óptimo interior finito porque el ingreso marginal es negativo.
    \item El monopolista fija un precio igual a cinco veces su costo marginal positivo.
    \item El índice de Lerner alcanza un valor superior a la unidad en el óptimo global.
    \item El monopolista maximiza beneficios produciendo en el punto de ingreso total mínimo.
\end{opts}

\q{Respecto al índice de Lerner $L = \frac{P - CMg}{P}$ en un monopolio no regulado con $CMg > 0$, señale la afirmación correcta:}
\begin{opts}
    \item Su valor es siempre mayor que la unidad cuando la demanda del mercado es elástica.
    \item Coincide con el inverso negativo de la elasticidad precio en el punto de equilibrio.
    \item Aumenta necesariamente cuando el costo fijo del monopolista se incrementa en el corto plazo.
    \item Permanece idéntico a cero siempre que la curva de demanda del mercado sea lineal y decreciente.
\end{opts}

\q{Un monopolista con costo marginal constante $c$ enfrenta una curva de demanda lineal $P = a - bQ$. Si el Estado impone un impuesto específico $t$ por unidad producida, ¿cuál es el impacto sobre el precio de equilibrio?}
\begin{opts}
    \item El precio aumenta exactamente en $t$, trasladándose el 100\% del gravamen a los consumidores.
    \item El precio aumenta en una proporción mayor que $t$ debido al poder de mercado del monopolista.
    \item El precio aumenta exactamente en $t/2$, repartiéndose la carga por igual entre ambas partes.
    \item El precio permanece inalterado porque el monopolista absorbe el impuesto con sus rentas.
\end{opts}

\q{¿Bajo qué condición un impuesto específico por unidad $t$ aplicado a un monopolista con costo marginal constante genera un sobretraslado ($\Delta P > t$) hacia el precio final?}
\begin{opts}
    \item Cuando la demanda es perfectamente inelástica en todo el rango relevante de precios.
    \item Cuando la demanda es lineal con intercepto vertical suficientemente elevado y pendiente baja.
    \item Cuando la función de demanda es estrictamente cóncava respecto al origen en el equilibrio.
    \item Cuando la demanda presenta elasticidad precio constante con $|\varepsilon| > 1$.
\end{opts}

\q{Un monopolista que opera dos plantas con funciones de costo $C_1(Q_1)$ y $C_2(Q_2)$ estrictamente convexas maximiza beneficios asignando producción positiva a ambas plantas cuando:}
\begin{opts}
    \item Iguala el costo marginal de cada planta con el ingreso marginal de la producción total.
    \item Iguala el costo medio de ambas plantas con el precio vigente en el mercado final.
    \item Asigna toda la producción a la planta que exhibe menor costo fijo medio en el corto plazo.
    \item Iguala el ingreso marginal de la planta más grande con el costo total medio de la menor.
\end{opts}

\q{Según la conjetura de Coase aplicada a un monopolista que vende un bien perfectamente duradero sin capacidad de compromiso intertemporal, ¿qué ocurre conforme el intervalo entre ajustes de precio tiende a cero?}
\begin{opts}
    \item El monopolista logra extraer la totalidad del excedente del consumidor mediante ventas sucesivas.
    \item El precio converge al costo marginal al competir el monopolista contra sus ventas futuras.
    \item El monopolista restringe la oferta indefinidamente y mantiene el precio estático de monopolio.
    \item Los consumidores adelantan todas sus compras al primer periodo pagando el precio de reserva máximo.
\end{opts}

\q{¿Cuál de las siguientes estrategias permite a un monopolista de bienes duraderos mitigar el problema de inconsistencia temporal planteado por Coase?}
\begin{opts}
    \item Aumentar la durabilidad física del producto para reducir la tasa de reposición en el mercado.
    \item Anunciar descuentos futuros frecuentes para incentivar la entrada de compradores con baja valoración.
    \item Alquilar el bien en lugar de venderlo o incluir cláusulas de nación más favorecida en sus contratos.
    \item Eliminar toda restricción de capacidad instalada para demostrar que puede abastecer al mercado completo.
\end{opts}

\q{En el análisis del costo social del monopolio, la diferencia fundamental entre el enfoque de Harberger y el enfoque de Tullock (búsqueda de rentas) radica en que Tullock sostiene que:}
\begin{opts}
    \item El triángulo de pérdida irrecuperable de eficiencia sobreestima gravemente el daño del monopolio.
    \item La transferencia de excedente del consumidor al monopolista siempre genera una ganancia neta social.
    \item El costo social se reduce a cero cuando el monopolista realiza discriminación de tercer grado.
    \item Los recursos reales disipados para obtener o defender rentas monopólicas también son pérdida social.
\end{opts}

\q{De acuerdo con el modelo de Spence sobre elección de calidad por un monopolista que cobra un precio uniforme, la provisión de calidad respecto al óptimo social depende de:}
\begin{opts}
    \item La valoración de la calidad por el consumidor marginal frente a la del consumidor promedio.
    \item La relación entre el costo fijo de producción y la elasticidad ingreso de la demanda agregada.
    \item El signo de la segunda derivada del costo total respecto a la cantidad producida en el corto plazo.
    \item La existencia de economías de alcance entre los mercados regulados y no regulados de la firma.
\end{opts}

\q{Para que una cantidad $Q^* > 0$ que satisface $IMg(Q^*) = CMg(Q^*)$ corresponda efectivamente a un máximo local de beneficios del monopolista, se requiere que en $Q^*$:}
\begin{opts}
    \item La pendiente del costo marginal sea estrictamente menor que la pendiente del ingreso marginal.
    \item La pendiente del ingreso marginal sea menor que la pendiente del costo marginal.
    \item La elasticidad precio de la demanda sea exactamente igual a menos uno en el equilibrio.
    \item El costo medio total sea estrictamente creciente e igual al precio de mercado vigente.
\end{opts}

\q{En un equilibrio con discriminación perfecta de precios (primer grado) sin costos de transacción, ¿cuál de los siguientes resultados se verifica simultáneamente?}
\begin{opts}
    \item El excedente del consumidor es máximo y el beneficio económico del monopolista es nulo.
    \item La producción es inferior a la competitiva pero el bienestar total alcanza su máximo teórico.
    \item La cantidad producida es socialmente eficiente y el monopolista captura todo el excedente.
    \item El precio cobrado a la última unidad vendida supera estrictamente al costo marginal de producción.
\end{opts}

\q{Un monopolista practica discriminación de tercer grado entre dos mercados separados A y B con idéntico costo marginal constante. Si en el equilibrio $|\varepsilon_A| < |\varepsilon_B|$, entonces:}
\begin{opts}
    \item El precio en el mercado B supera al del mercado A porque su demanda es más sensible al precio.
    \item El ingreso marginal en el mercado A supera estrictamente al ingreso marginal en el mercado B.
    \item El monopolista cierra el mercado A para concentrar toda su producción en el mercado elástico.
    \item El precio es mayor en el mercado A porque su demanda es relativamente menos elástica.
\end{opts}

\q{Según el teorema de Schmalensee y Varian sobre el bienestar social bajo discriminación de tercer grado respecto al precio único (con costo marginal constante), una condición necesaria para que el bienestar total aumente es que:}
\begin{opts}
    \item La cantidad total vendida en el conjunto de los mercados aumente estrictamente.
    \item Los precios disminuyan simultáneamente en todos los segmentos de consumidores atendidos.
    \item Todas las demandas segmentadas sean lineales y se atiendan ambos mercados bajo precio único.
    \item El excedente del consumidor en el mercado menos elástico aumente tras la discriminación.
\end{opts}

\q{Si un monopolista con costo marginal constante atiende dos mercados con demandas lineales que ya estaban ambos activos bajo precio único, pasar a discriminación de tercer grado provoca que:}
\begin{opts}
    \item La producción agregada aumente y el bienestar social mejore de forma inequívoca.
    \item La producción agregada permanezca constante y el bienestar social disminuya.
    \item La producción agregada disminuya y el beneficio del monopolista también se reduzca.
    \item Los precios en ambos mercados converjan al costo marginal de largo plazo.
\end{opts}

\q{En el diseño óptimo de un menú de contratos bajo discriminación de precios de segundo grado (dos tipos de consumidores: alta y baja valoración), la propiedad conocida como «no distorsión en el tope» establece que:}
\begin{opts}
    \item El consumidor de baja valoración recibe la cantidad socialmente eficiente sin renta informacional.
    \item Ambos tipos de consumidores reciben cantidades estrictamente inferiores a las socialmente eficientes.
    \item El consumidor de alta valoración recibe la cantidad eficiente y obtiene una renta informacional positiva.
    \item El monopolista extrae todo el excedente del consumidor de alta valoración fijando un precio marginal alto.
\end{opts}

\q{En el problema estándar de discriminación de segundo grado con dos tipos de consumidores (bajo $L$ y alto $H$), ¿cuáles son las dos restricciones que resultan activas (con igualdad) en el contrato óptimo del monopolista?}
\begin{opts}
    \item La restricción de participación del tipo $H$ y la compatibilidad de incentivos del tipo $L$.
    \item Ambas restricciones de participación, dejando inactivas las de compatibilidad de incentivos.
    \item Ambas restricciones de compatibilidad de incentivos, garantizando excedente positivo al tipo $L$.
    \item La restricción de participación del tipo $L$ y la compatibilidad de incentivos del tipo $H$.
\end{opts}

\q{¿Cuál de las siguientes condiciones es indispensable para que una empresa pueda sostener una estrategia rentable de discriminación de precios de tercer grado?}
\begin{opts}
    \item Poder de mercado, segmentación de demandas distintas y ausencia de arbitraje entre grupos.
    \item Retornos crecientes a escala en todas las plantas y costos marginales nulos en el corto plazo.
    \item Información perfecta sobre la disposición a pagar individual de cada consumidor del mercado.
    \item Ausencia total de bienes sustitutos cercanos y regulación estatal de precios máximos.
\end{opts}

\q{En el sector saneamiento regulado por SUNASS, la aplicación de tarifas por bloques crecientes (IBT) en el consumo residencial de agua potable busca principalmente:}
\begin{opts}
    \item Maximizar los beneficios monopólicos de la EPS explotando la inexistencia de micromedición.
    \item Conciliar equidad en consumos básicos con señales de conservación en altos consumos.
    \item Igualar el precio marginal de todos los bloques al costo medio histórico de la infraestructura.
    \item Eliminar la necesidad de aplicar un cargo fijo mensual para cubrir los costos de facturación.
\end{opts}

\q{En discriminación de segundo grado mediante versionado (por ejemplo, software estándar vs. profesional), ¿por qué un monopolista puede incurrir en costos adicionales para degradar deliberadamente la versión básica?}
\begin{opts}
    \item Para reducir los costos marginales de producción de la versión de alta calidad mediante economías de escala.
    \item Para cumplir con los requisitos técnicos de compatibilidad exigidos por los organismos reguladores.
    \item Para relajar la restricción de compatibilidad de incentivos del consumidor de alta valoración.
    \item Para asegurar que los consumidores de baja valoración obtengan una renta informacional positiva.
\end{opts}

\q{Cuando un monopolista fija un precio elevado al lanzar un producto nuevo y lo reduce meses después para captar consumidores más sensibles al precio, practica discriminación intertemporal. Esta estrategia fracasa si:}
\begin{opts}
    \item Los consumidores de alta valoración tienen una tasa de descuento subjetiva muy elevada.
    \item Los costos marginales de producción disminuyen progresivamente por efecto de la curva de aprendizaje.
    \item El producto es altamente perecedero y no puede ser almacenado por intermediarios comerciales.
    \item Los consumidores de alta valoración son muy pacientes y anticipan perfectamente la rebaja futura.
\end{opts}

\q{Una diferencia clave entre la discriminación de precios de segundo grado (tarifas no lineales) y la de tercer grado radica en que en la discriminación de segundo grado:}
\begin{opts}
    \item El vendedor no observa el tipo de cada cliente y depende de la autoselección mediante el menú.
    \item El vendedor cobra siempre un precio unitario constante dentro de cada segmento observable.
    \item El bienestar social es necesariamente inferior al que se obtiene bajo un precio uniforme.
    \item El arbitraje entre consumidores resulta imposible incluso cuando los bienes son fácilmente revendibles.
\end{opts}

\q{Si una tarifa no lineal presenta descuentos por volumen (tarifa por bloques decrecientes), el precio medio pagado por unidad consumida es:}
\begin{opts}
    \item Estrictamente inferior al precio marginal del último bloque consumido por el usuario.
    \item Decreciente con la cantidad consumida y mayor al precio marginal del último bloque.
    \item Constante para cualquier nivel de consumo siempre que no exista un cargo fijo de conexión.
    \item Igual al costo marginal de corto plazo en todos los tramos de la escala tarifaria.
\end{opts}

\q{¿En qué situación la discriminación de precios de tercer grado genera inequívocamente un aumento del bienestar social respecto a un precio uniforme?}
\begin{opts}
    \item Cuando la elasticidad precio es idéntica en todos los mercados y el costo marginal es creciente.
    \item Cuando ambos mercados son atendidos bajo precio único y sus demandas son estrictamente cóncavas.
    \item Cuando permite atender a un mercado de baja valoración que quedaría excluido con precio único.
    \item Cuando el monopolista reduce su producción en el mercado elástico para elevar el precio medio.
\end{opts}

\q{El uso de cupones de descuento que requieren esfuerzo o tiempo de recorte por parte del comprador constituye un mecanismo económico de:}
\begin{opts}
    \item Discriminación de primer grado basada en información perfecta sobre la renta del consumidor.
    \item Competencia a la Bertrand diseñada para eliminar las rentas económicas en el corto plazo.
    \item Integración vertical hacia adelante para evitar el problema de la doble marginalización.
    \item Autoselección que separa clientes según su costo del tiempo y elasticidad precio.
\end{opts}

\q{Si un monopolista que discriminaba en tercer grado entre dos mercados pierde súbitamente la capacidad de impedir el arbitraje entre ellos, el precio único resultante se ubicará:}
\begin{opts}
    \item En un nivel intermedio entre el precio del mercado menos elástico y el del más elástico.
    \item Por encima del precio que anteriormente regía en el mercado de demanda menos elástica.
    \item Por debajo del precio que anteriormente regía en el mercado de demanda más elástica.
    \item Exactamente en el nivel del costo marginal de producción por efecto de la reventa libre.
\end{opts}

\q{Un monopolista enfrenta $N$ consumidores idénticos con demanda individual $q(p)$ y produce con costo marginal constante $c$. Según Walter Oi, la tarifa en dos partes óptima $(T, p)$ consiste en fijar:}
\begin{opts}
    \item Un precio por uso igual al precio de monopolio puro y un cargo fijo igual a cero.
    \item Un precio por uso igual a $c$ y un cargo fijo igual al excedente individual en $c$.
    \item Un precio por uso por debajo del costo marginal y un cargo fijo igual al costo fijo medio.
    \item Un precio por uso igual al costo medio y un cargo fijo que cubra la pérdida de eficiencia.
\end{opts}

\q{Si un monopolista aplica una única tarifa en dos partes $(T, p)$ a dos grupos de consumidores heterogéneos (alta y baja demanda) y decide atender a ambos grupos, el precio óptimo por uso $p^*$ satisfará:}
\begin{opts}
    \item $p^* = CMg$ para maximizar el tamaño del excedente total en ambos grupos simultáneamente.
    \item $p^* < CMg$ con el objetivo de incentivar el consumo del grupo de baja demanda.
    \item $p^* > CMg$ para extraer más excedente del grupo alto sin expulsar al grupo bajo.
    \item $p^* = CMe$ independientemente de la proporción relativa de cada tipo de consumidor.
\end{opts}

\q{En el análisis clásico de Stigler y Adams-Yellen sobre empaquetamiento puro (pure bundling) de dos bienes con costo marginal nulo, esta estrategia incrementa los beneficios del monopolista principalmente cuando:}
\begin{opts}
    \item Las valoraciones de los consumidores por ambos bienes están perfectamente correlacionadas de forma positiva.
    \item Todos los consumidores poseen exactamente la misma disposición a pagar por cada bien por separado.
    \item El costo marginal de producción de ambos bienes supera la valoración máxima de los compradores.
    \item Las valoraciones de los consumidores por los dos bienes están negativamente correlacionadas.
\end{opts}

\q{¿Por qué el empaquetamiento mixto (ofrecer los bienes tanto en paquete como por separado) puede dominar estrictamente al empaquetamiento puro cuando existen costos marginales positivos?}
\begin{opts}
    \item Evita vender un bien a consumidores cuya valoración por dicho bien es inferior a su costo marginal.
    \item Elimina por completo el costo fijo de distribución al estandarizar las presentaciones del producto.
    \item Permite cobrar por el paquete un precio estrictamente superior a la suma de los precios individuales.
    \item Impide la entrada de competidores al garantizar que todos los clientes compren el paquete completo.
\end{opts}

\q{Cuando un fabricante de impresoras vende la máquina a un precio cercano al costo y obliga a comprar exclusivamente sus propios cartuchos de tinta a un precio elevado, esta venta atada actúa principalmente como:}
\begin{opts}
    \item Una estrategia de precios predatorios para expulsar a los productores de papel del mercado.
    \item Un mecanismo de medición (metering) que permite discriminar según la intensidad de uso.
    \item Una forma de regulación voluntaria que elimina la pérdida de eficiencia asignativa.
    \item Un esquema de subsidio cruzado hacia los usuarios corporativos de alto volumen de impresión.
\end{opts}

\q{Según la crítica de la Escuela de Chicago a la teoría del apalancamiento mediante venta atada, si un monopolista del bien A lo ata a un bien complementario B producido competitivamente en proporciones fijas:}
\begin{opts}
    \item Duplica automáticamente sus beneficios monopolizando el mercado del bien complementario B.
    \item Reduce el precio del bien A por debajo de su costo marginal para expandir el mercado de B.
    \item No puede aumentar su beneficio total porque ya capturaba toda la renta en el precio de A.
    \item Genera una pérdida de eficiencia adicional equivalente al doble del triángulo de Harberger.
\end{opts}

\q{En el modelo clásico de fijación de precios pico (Steiner-Boiteux) con capacidad rígida, costo operativo $b$ y costo marginal de capacidad $\beta$, en el caso de pico firme (demanda valle claramente inferior a la capacidad óptima), los precios eficientes son:}
\begin{opts}
    \item $P_{\text{pico}} = b + \beta/2$ y $P_{\text{valle}} = b + \beta/2$, dividiendo la capacidad en partes iguales.
    \item $P_{\text{pico}} = b$ y $P_{\text{valle}} = b + \beta$, financiando la planta con la demanda más inelástica.
    \item $P_{\text{pico}} = b + \beta$ y $P_{\text{valle}} = b + \beta$, exigiendo cobertura total en ambos periodos.
    \item $P_{\text{pico}} = b + \beta$ y $P_{\text{valle}} = b$, asignando el costo de capacidad al periodo pico.
\end{opts}

\q{En el modelo de precios pico, ¿qué sucede en el caso de «pico cambiante» (cuando al fijar $P_{\text{valle}} = b$ la demanda del periodo valle superaría la capacidad instalada)?}
\begin{opts}
    \item Ambos periodos operan a plena capacidad y la suma de sus precios equivale a $2b + \beta$.
    \item Se fija $P_{\text{valle}} = b$ y se raciona el exceso de demanda mediante cortes programados.
    \item El costo de capacidad $\beta$ recae íntegramente sobre el periodo valle y cero en el pico.
    \item La capacidad óptima se expande hasta que el precio en ambos periodos sea igual a cero.
\end{opts}

\q{¿Cuál es el principal beneficio asignativo de reemplazar una tarifa uniforme por un esquema de precios pico en infraestructura de agua potable o electricidad?}
\begin{opts}
    \item Incrementa la variabilidad horaria del consumo para aprovechar mejor las economías de alcance.
    \item Reduce la necesidad de sobreinvertir en capacidad ociosa y suaviza el perfil de demanda.
    \item Aumenta el consumo total en las horas de máxima congestión al disminuir el precio medio.
    \item Permite al regulador prescindir de la medición del consumo horario o estacional de los usuarios.
\end{opts}

\q{En comparación con un monopolista que cobra un precio lineal uniforme, un monopolista que implementa una tarifa en dos partes óptima con consumidores homogéneos logra:}
\begin{opts}
    \item Aumentar el excedente del consumidor a costa de reducir sus propios beneficios económicos.
    \item Mantener inalterada la pérdida de eficiencia social mientras transfiere rentas al Estado.
    \item Eliminar la pérdida irrecuperable de eficiencia asignativa incrementando sus beneficios.
    \item Reducir la cantidad producida para elevar simultáneamente el cargo fijo y el precio marginal.
\end{opts}

\q{Cuando un monopolista enfrenta consumidores de baja y alta demanda y solo puede ofrecer una única tarifa en dos partes $(T, p)$, ¿cuándo le conviene excluir del mercado a los consumidores de baja demanda?}
\begin{opts}
    \item Siempre que el costo marginal de producción sea estrictamente positivo en el corto plazo.
    \item Cuando la proporción de consumidores de baja demanda es mucho mayor que la de alta demanda.
    \item Cuando la diferencia entre la disposición a pagar de ambos grupos es muy pequeña.
    \item Cuando el grupo alto es numeroso y su excedente supera ampliamente al del bajo.
\end{opts}

\q{En el diseño de tarifas de agua potable en dos partes (cargo fijo + tarifa volumétrica), fijar el cargo fijo muy por encima de los costos fijos de facturación y medición puede generar un problema de:}
\begin{opts}
    \item Exclusión o morosidad en hogares vulnerables, afectando el objetivo de acceso universal.
    \item Sobreconsumo excesivo de agua debido a un precio marginal artificialmente superior al costo.
    \item Arbitraje masivo de agua potable entre usuarios residenciales y usuarios industriales.
    \item Incompatibilidad técnica con el principio de subaditividad de costos en las redes de distribución.
\end{opts}

\q{Según el modelo de Whinston (1990) sobre venta atada estratégica, ¿en qué circunstancia el empaquetamiento puro puede utilizarse racionalmente para excluir a un rival en un mercado complementario?}
\begin{opts}
    \item Cuando el mercado complementario opera con retornos constantes a escala y libre entrada sin costos fijos.
    \item Cuando el mercado complementario tiene economías de escala y el atado compromete una respuesta agresiva.
    \item Cuando el monopolista carece de poder de mercado en el bien principal pero tiene costos marginales bajos.
    \item Cuando los consumidores valoran ambos bienes de forma independiente y no existen costos hundidos.
\end{opts}

\q{La condición necesaria y suficiente para que una industria sea un monopolio natural para un vector de producción $Q$ es que su función de costos sea:}
\begin{opts}
    \item Estrictamente decreciente en el costo medio para todos los niveles posibles de producción.
    \item Homogénea de grado mayor que uno en todos los factores productivos empleados por la firma.
    \item Estrictamente subaditiva en $Q$: una sola firma produce $Q$ a menor costo que varias.
    \item Lineal en los costos variables con ausencia total de costos hundidos en la infraestructura.
\end{opts}

\q{En una industria uniproducto con curva de costo medio en forma de U, ¿puede existir monopolio natural en un nivel de producción donde el costo medio ya es creciente?}
\begin{opts}
    \item No, porque las economías de escala globales son condición necesaria y suficiente de subaditividad.
    \item No, porque cuando el costo medio crece, el costo marginal es inferior al costo medio total.
    \item Sí, pero únicamente si la demanda del mercado es completamente inelástica en ese tramo.
    \item Sí, porque la subaditividad puede persistir en un rango más allá del mínimo del costo medio.
\end{opts}

\q{En una empresa prestadora de servicios de saneamiento (EPS) que provee conjuntamente agua potable y alcantarillado, la subaditividad multiproducto requiere la interacción de:}
\begin{opts}
    \item Economías de escala específicas (o radiales) combinadas con economías de alcance.
    \item Retornos decrecientes a escala en cada servicio y costos marginales crecientes en la red.
    \item Separación contable absoluta y ausencia total de insumos compartidos entre ambos servicios.
    \item Elasticidades cruzadas negativas entre la demanda de agua potable y la de alcantarillado.
\end{opts}

\q{¿A qué se refiere el fenómeno de «descreme del mercado» (cream-skimming) en un monopolio natural regulado con precios uniformes o subsidios cruzados?}
\begin{opts}
    \item A la estrategia del incumbente de reducir la calidad en los segmentos de menores ingresos.
    \item Al ingreso de rivales solo en segmentos rentables, afectando la viabilidad financiera.
    \item A la fijación de tarifas por debajo del costo incremental para expulsar a rivales eficientes.
    \item A la captura del regulador por parte de las empresas prestadoras para inflar la base de capital.
\end{opts}

\q{Según la teoría de los mercados perfectamente contestables, un monopolio natural se ve disciplinado a fijar precios iguales al costo medio (beneficio económico nulo) incluso sin competidores actuales cuando:}
\begin{opts}
    \item El regulador impone un techo de precios basado en la tasa de retorno histórica de la industria.
    \item Existen barreras legales a la entrada pero el bien cuenta con sustitutos importados cercanos.
    \item La entrada es libre y la salida no presenta costos hundidos, permitiendo entrada relámpago.
    \item La función de costos exhibe economías de alcance pero deseconomías de escala pronunciadas.
\end{opts}

\q{El dilema fundamental al regular un monopolio natural con economías de escala globales ($CMg < CMe$) consiste en que:}
\begin{opts}
    \item Fijar $P = CMe$ genera pérdidas financieras permanentes que requieren subsidios fiscales.
    \item Fijar $P = CMg$ genera beneficios extraordinarios que incentivan el ingreso ineficiente.
    \item Cualquier tarifa lineal regulada elimina los costos hundidos de la red de distribución.
    \item Fijar $P = CMg$ logra eficiencia asignativa pero genera déficit al monopolista.
\end{opts}

\q{En un monopolio natural multiproducto con demandas independientes y restricción de autofinanciamiento ($\pi = 0$), la regla de precios óptimos de segundo mejor de Ramsey-Boiteux establece que:}
\begin{opts}
    \item El margen relativo sobre el costo marginal es inversamente proporcional a $|\varepsilon_i|$.
    \item Todos los servicios deben recibir exactamente el mismo margen porcentual sobre su costo marginal.
    \item Los servicios con demanda más elástica deben soportar una mayor proporción de los costos comunes.
    \item El precio de cada servicio debe igualarse estrictamente a su costo marginal de largo plazo.
\end{opts}

\q{¿Cuál es la principal crítica desde el punto de vista de la equidad social a la aplicación pura de los precios de Ramsey-Boiteux en servicios básicos como el agua potable?}
\begin{opts}
    \item Genera un déficit financiero crónico que impide la reposición de la infraestructura hídrica.
    \item Impone mayores márgenes sobre bienes esenciales inelásticos consumidos por hogares pobres.
    \item Induce una sobrecapitalización sistemática en la expansión de redes de alcantarillado.
    \item Maximiza la pérdida irrecuperable de eficiencia al distorsionar las cantidades consumidas.
\end{opts}

\q{En la fórmula de Ramsey-Boiteux $\frac{P_i - CMg_i}{P_i} = \frac{R}{|\varepsilon_i|}$, donde $R = \frac{\lambda}{1+\lambda}$ es el número de Ramsey, ¿qué ocurre a medida que $R$ se aproxima a $1$?}
\begin{opts}
    \item Los precios regulados convergen hacia los costos marginales de primer mejor ($P_i = CMg_i$).
    \item La restricción de autofinanciamiento se vuelve no activa y el monopolio obtiene superávit.
    \item Los precios de Ramsey convergen hacia los precios que fijaría un monopolista no regulado.
    \item Las elasticidades precio de todos los servicios se igualan automáticamente a la unidad.
\end{opts}

\q{Según el criterio de Faulhaber para una empresa multiproducto que cubre sus costos totales ($R_1 + R_2 = C(q_1, q_2)$), una estructura tarifaria está libre de subsidios cruzados si y solo si los ingresos de cada servicio $i$ satisfacen:}
\begin{opts}
    \item Son mayores que su costo individual aislado (stand-alone) y menores que su costo incremental.
    \item Son exactamente iguales al costo marginal multiplicado por la cantidad producida de cada bien.
    \item Superan el costo total conjunto dividido en partes iguales entre el número de servicios.
    \item Se ubican entre su costo incremental y su costo de provisión independiente (stand-alone).
\end{opts}

\q{Si una empresa regulada cobra a un segmento industrial un ingreso total superior al costo de provisión independiente (stand-alone) de dicho servicio, el criterio de Faulhaber indica que:}
\begin{opts}
    \item El segmento industrial está otorgando un subsidio cruzado a otros segmentos de la empresa.
    \item El segmento industrial está recibiendo un subsidio cruzado proveniente de los usuarios residenciales.
    \item La empresa opera con deseconomías de alcance y debe escindirse inmediatamente en dos firmas.
    \item La estructura tarifaria coincide exactamente con el óptimo de primer mejor de eficiencia asignativa.
\end{opts}

\q{Cuando el regulador fija un precio máximo vinculante $P_{\max}$ por debajo del precio de monopolio pero por encima del costo marginal competitivo ($CMg < P_{\max} < P_m$), ¿qué ocurre con la cantidad ofrecida por el monopolista?}
\begin{opts}
    \item Disminuye estrictamente porque el menor precio reduce el incentivo marginal a producir.
    \item Aumenta porque el ingreso marginal se vuelve constante e igual a $P_{\max}$ hasta la demanda.
    \item Permanece inalterada mientras el costo fijo medio no supere al nuevo precio regulado.
    \item Se iguala inmediatamente a la cantidad competitiva sin importar cuán cerca esté $P_{\max}$ de $P_m$.
\end{opts}

\q{Se dice que una función de costos multiproducto exhibe economías de alcance (scope) entre los bienes 1 y 2 cuando:}
\begin{opts}
    \item El costo marginal de cada bien disminuye al incrementar su propia escala de producción.
    \item El costo total conjunto es estrictamente mayor que la suma de los costos de producción separada.
    \item Producir ambos bienes conjuntamente cuesta menos que producirlos por separado.
    \item La elasticidad costo-producto radial es estrictamente mayor que la unidad en todo el dominio.
\end{opts}

\q{En industrias de red verticalmente integradas donde un competidor solicita acceso a una instalación esencial del incumbente, la regla del componente eficiente (ECPR o regla Baumol-Willig) establece que el precio de acceso debe igualar:}
\begin{opts}
    \item El costo marginal de corto plazo del acceso, excluyendo cualquier costo de oportunidad del incumbente.
    \item El precio minorista final del incumbente multiplicado por el índice de Lerner del mercado aguas abajo.
    \item El costo fijo medio histórico de la red dividido entre el número total de competidores entrantes.
    \item El costo directo del acceso más el costo de oportunidad (margen perdido) del incumbente.
\end{opts}

\q{Según el modelo de Averch-Johnson, cuando un monopolio regulado enfrenta un techo a la tasa de retorno sobre el capital $s$ estrictamente mayor que el costo de mercado del capital $r$ ($s > r$) pero inferior a la tasa no regulada, la empresa:}
\begin{opts}
    \item Elige una relación capital-trabajo superior a la que minimiza costos para ese nivel de producto.
    \item Sustituye capital por trabajo de manera excesiva para reducir su base tarifaria regulada.
    \item Minimiza sus costos productivos de forma eficiente pero produce una cantidad inferior a la monopólica.
    \item Alcanza la eficiencia asignativa ($P = CMg$) al reducir automáticamente los gastos operativos.
\end{opts}

\q{¿Cuál es la principal ventaja teórica de la regulación por precios tope (price cap o $RPI - X$) frente a la regulación tradicional por tasa de retorno?}
\begin{opts}
    \item Elimina por completo el riesgo financiero de la empresa regulada ante choques de demanda.
    \item Genera incentivos de alto poder para bajar costos al retener ahorros en el periodo.
    \item Garantiza que los beneficios extraordinarios de la firma sean nulos en cada año regulatorio.
    \item Asegura un nivel óptimo de calidad del servicio sin necesidad de fiscalización adicional.
\end{opts}

\q{Entre las principales limitaciones prácticas de un esquema puro de regulación por precios tope ($RPI - X$) se encuentran:}
\begin{opts}
    \item La tendencia a la sobrecapitalización excesiva y la ausencia de incentivos para reducir OPEX.
    \item La imposibilidad de ajustar tarifas ante procesos inflacionarios de la economía nacional.
    \item El riesgo de degradar calidad y el efecto trinquete al recalibrar el factor $X$.
    \item La eliminación automática de la asimetría de información entre el regulador y la empresa.
\end{opts}

\q{En el mecanismo de competencia por comparación (yardstick competition) propuesto por Shleifer (1985) para regular monopolios naturales regionales como las EPS, la tarifa permitida a una empresa depende de:}
\begin{opts}
    \item Sus propios costos contables históricos más un margen fijo aprobado por el directorio.
    \item El precio promedio observado en mercados competitivos internacionales de bienes transables.
    \item La disposición a pagar declarada voluntariamente por los usuarios en audiencias públicas.
    \item Los costos observados de empresas comparables (ajustados), excluyendo los propios.
\end{opts}

\q{En el mecanismo regulatorio de Loeb y Magat (1979) con demanda conocida pero costos desconocidos por el regulador, se induce al monopolista a producir la cantidad eficiente ($P = CMg$) mediante:}
\begin{opts}
    \item Una transferencia al monopolista igual al excedente del consumidor en el precio elegido.
    \item La imposición de un impuesto de suma alzada equivalente al costo fijo declarado por la firma.
    \item La licitación de la concesión a la empresa que ofrezca la mayor tasa de retorno sobre el capital.
    \item La fijación administrativa de un precio igual al costo medio de una empresa modelo hipotética.
\end{opts}

\q{En el modelo principal-agente de regulación con información asimétrica de Laffont y Tirole, existe un dilema fundamental (trade-off) entre:}
\begin{opts}
    \item La subaditividad de la red de distribución y la elasticidad ingreso de los usuarios.
    \item La extracción de rentas informacionales y los incentivos al esfuerzo reductor de costos.
    \item El tamaño del costo hundido y la tasa de depreciación contable de los activos regulados.
    \item La competencia a la Cournot aguas arriba y la discriminación de tercer grado aguas abajo.
\end{opts}

\q{En el esquema óptimo de Laffont-Tirole con costo social de los fondos públicos positivo, el regulador ofrece un menú de contratos donde la empresa más eficiente selecciona:}
\begin{opts}
    \item Un contrato de reembolso total de costos (cost-plus) sin ninguna renta informacional.
    \item Un contrato con alta participación del regulador en los sobrecostos y esfuerzo distorsionado.
    \item Un contrato de precio fijo con esfuerzo eficiente y renta informacional positiva.
    \item Un contrato de tasa de retorno fija que induce sobrecapitalización para revelar su tipo.
\end{opts}

\q{La propuesta de Harold Demsetz de «competencia por el mercado» como alternativa a la regulación tradicional de monopolios naturales consiste en:}
\begin{opts}
    \item Dividir físicamente las redes de distribución entre múltiples operadores simultáneos en la misma zona.
    \item Subastar la franquicia al postor que ofrezca el mayor pago de suma alzada al Tesoro Público.
    \item Permitir que el monopolista fije libremente precios a cambio de publicar sus estados financieros.
    \item Licitar la concesión a quien ofrezca operar cobrando la tarifa más baja al usuario.
\end{opts}

\q{En el marco regulatorio peruano de los servicios de saneamiento supervisado por SUNASS, la determinación de la fórmula tarifaria quinquenal de una EPS se sustenta en:}
\begin{opts}
    \item Un Plan Maestro Optimizado (PMO) y el equilibrio financiero a valor presente con la tasa WACC.
    \item La igualación anual de la tarifa media al costo marginal de corto plazo sin considerar inversiones.
    \item Una subasta mensual de volúmenes de agua cruda entre los usuarios comerciales e industriales.
    \item La aplicación exclusiva de precios de Ramsey puros sin considerar subsidios cruzados sociales.
\end{opts}

\q{Dentro del esquema regulatorio aplicado por SUNASS en el Perú, ¿con qué finalidad se incorpora el reajuste automático de tarifas cuando el Índice de Precios al por Mayor (IPM) acumula una variación determinada?}
\begin{opts}
    \item Para penalizar a la EPS cuando incumple las metas de continuidad y presión del servicio.
    \item Para preservar el valor real de los ingresos tarifarios frente a la inflación.
    \item Para trasladar inmediatamente a los usuarios cualquier ineficiencia administrativa de la EPS.
    \item Para sustituir la necesidad de evaluar el factor de productividad $X$ en la revisión tarifaria.
\end{opts}

\q{En la política tarifaria de SUNASS, la transición desde subsidios cruzados generalizados hacia subsidios cruzados focalizados (usando información socioeconómica como el SISFOH o planos estratificados) busca:}
\begin{opts}
    \item Eliminar toda diferenciación entre usuarios domésticos, comerciales, industriales y estatales.
    \item Asegurar que las tarifas de los bloques industriales sean inferiores al costo medio de operación.
    \item Reducir los errores de inclusión y exclusión, dirigiendo el subsidio a hogares vulnerables.
    \item Reemplazar por completo los recursos del presupuesto de la EPS por transferencias directas del MEF.
\end{opts}

\subsection*{Parte II. Aplicaciones prácticas y ejercicios cuantitativos (154 a 180)}

\q{Un monopolista enfrenta la función de demanda inversa $P = 120 - 2Q$ y opera con una función de costo total $CT(Q) = 20Q + 100$. ¿Cuáles son el precio óptimo $P^*$, la cantidad óptima $Q^*$ y el beneficio máximo $\pi^*$?}
\begin{opts}
    \item $P^* = 60$, $Q^* = 30$, $\pi^* = 1100$
    \item $P^* = 80$, $Q^* = 20$, $\pi^* = 1100$
    \item $P^* = 70$, $Q^* = 25$, $\pi^* = 1250$
    \item $P^* = 70$, $Q^* = 25$, $\pi^* = 1150$
\end{opts}

\q{En un mercado con demanda $P = 100 - 2Q$ abastecido por un monopolista con costo marginal constante $CMg = 20$, ¿a cuánto asciende la pérdida irrecuperable de eficiencia (triángulo de Harberger) respecto a la competencia perfecta?}
\begin{opts}
    \item $400$ unidades monetarias
    \item $800$ unidades monetarias
    \item $200$ unidades monetarias
    \item $600$ unidades monetarias
\end{opts}

\q{Un monopolista maximizador de beneficios registra en equilibrio un índice de Lerner igual a $L = 0.40$ y un costo marginal de $CMg = 30$. ¿Cuáles son el precio de equilibrio $P^*$ y la elasticidad precio $\varepsilon$ en ese punto?}
\begin{opts}
    \item $P^* = 42.0$ y $\varepsilon = -1.67$
    \item $P^* = 50.0$ y $\varepsilon = -2.50$
    \item $P^* = 75.0$ y $\varepsilon = -0.40$
    \item $P^* = 50.0$ y $\varepsilon = -1.50$
\end{opts}

\q{Un monopolista con costo marginal constante $CMg = 18$ enfrenta una curva de demanda de elasticidad precio constante igual a $\varepsilon = -2.5$. ¿Cuál es el precio óptimo que maximiza sus beneficios?}
\begin{opts}
    \item $P^* = 22.50$
    \item $P^* = 45.00$
    \item $P^* = 30.00$
    \item $P^* = 27.00$
\end{opts}

\q{Un monopolista enfrenta la demanda $P = 100 - Q$ y tiene $CMg = 20$. Si el Estado introduce un impuesto específico de $t = 12$ por unidad, ¿cuál será el nuevo precio de equilibrio $P_1^*$ y el incremento de precio $\Delta P$?}
\begin{opts}
    \item $P_1^* = 72$ y $\Delta P = 12$
    \item $P_1^* = 64$ y $\Delta P = 4$
    \item $P_1^* = 68$ y $\Delta P = 8$
    \item $P_1^* = 66$ y $\Delta P = 6$
\end{opts}

\q{Un monopolista con $CMg = 20$ enfrenta una demanda isoelástica con $\varepsilon = -3$. Si se aplica un impuesto específico de $t = 6$ por unidad vendida, ¿en cuánto aumenta el precio de equilibrio?}
\begin{opts}
    \item $\Delta P = 9.00$
    \item $\Delta P = 6.00$
    \item $\Delta P = 3.00$
    \item $\Delta P = 4.50$
\end{opts}

\q{Un monopolista enfrenta la demanda $P = 100 - Q$ y produce en dos plantas con funciones de costo $C_1(Q_1) = Q_1^2$ y $C_2(Q_2) = 2Q_2^2$. ¿Cuál es la asignación óptima $(Q_1^*, Q_2^*)$ y el precio de equilibrio $P^*$?}
\begin{opts}
    \item $Q_1^* = 15$, $Q_2^* = 15$, $P^* = 70$
    \item $Q_1^* = 20$, $Q_2^* = 10$, $P^* = 70$
    \item $Q_1^* = 10$, $Q_2^* = 20$, $P^* = 70$
    \item $Q_1^* = 24$, $Q_2^* = 12$, $P^* = 64$
\end{opts}

\q{Un monopolista enfrenta la demanda $P = 100 - Q$ y su función de costo total es $CT(Q) = 0.5Q^2 + 10Q + 100$. ¿Cuál es el beneficio económico máximo que puede obtener?}
\begin{opts}
    \item $\pi^* = 1350$
    \item $\pi^* = 1150$
    \item $\pi^* = 1250$
    \item $\pi^* = 1450$
\end{opts}

\q{Un monopolista con $CT(Q) = 20Q + 100$ enfrenta la demanda $P = 80 - 2Q$ y logra aplicar discriminación perfecta de precios (primer grado). ¿Cuál es el beneficio máximo $\pi^*$ que obtiene?}
\begin{opts}
    \item $\pi^* = 450$
    \item $\pi^* = 900$
    \item $\pi^* = 350$
    \item $\pi^* = 800$
\end{opts}

\q{Un monopolista con costo marginal constante $CMg = 20$ atiende dos mercados separados con demandas $P_1 = 100 - Q_1$ y $P_2 = 60 - Q_2$. Si aplica discriminación de tercer grado, ¿cuáles son los precios óptimos $P_1^*$ y $P_2^*$?}
\begin{opts}
    \item $P_1^* = 60$ y $P_2^* = 40$
    \item $P_1^* = 50$ y $P_2^* = 50$
    \item $P_1^* = 70$ y $P_2^* = 30$
    \item $P_1^* = 80$ y $P_2^* = 40$
\end{opts}

\q{Una empresa discrimina precios entre dos grupos de clientes con elasticidades constantes $\varepsilon_1 = -2$ y $\varepsilon_2 = -4$, siendo su costo marginal igual a $CMg = 15$ en ambos segmentos. ¿Cuál es la razón de precios óptimos $P_1^* / P_2^*$?}
\begin{opts}
    \item $P_1^* / P_2^* = 2.00$
    \item $P_1^* / P_2^* = 1.50$
    \item $P_1^* / P_2^* = 1.25$
    \item $P_1^* / P_2^* = 0.67$
\end{opts}

\q{Con demandas separadas $P_1 = 100 - Q_1$ y $P_2 = 60 - Q_2$, costo marginal $CMg = 20$ y costo fijo nulo, ¿en cuánto aumenta el beneficio del monopolista al pasar de un precio único a discriminación de tercer grado?}
\begin{opts}
    \item $400$ unidades monetarias
    \item $100$ unidades monetarias
    \item $200$ unidades monetarias
    \item $300$ unidades monetarias
\end{opts}

\q{Un club monopólico atiende consumidores idénticos, cada uno con demanda inversa $p = 50 - 0.5q$, y opera con un costo marginal constante $CMg = 10$. ¿Cuál es la tarifa en dos partes óptima $(T^*, p^*)$?}
\begin{opts}
    \item $T^* = 800$ y $p^* = 30$
    \item $T^* = 1200$ y $p^* = 20$
    \item $T^* = 2000$ y $p^* = 10$
    \item $T^* = 1600$ y $p^* = 10$
\end{opts}

\q{Un monopolista con $CMg = 4$ atiende a $1$ consumidor de alta demanda ($q_H = 18 - p$) y $2$ consumidores de baja demanda ($q_L = 12 - p$) mediante una única tarifa en dos partes $(T, p)$. ¿Cuál es el precio óptimo $p^*$ y cargo fijo $T^*$ que maximizan su beneficio atendiendo a todos?}
\begin{opts}
    \item $p^* = 6$ y $T^* = 18$
    \item $p^* = 4$ y $T^* = 32$
    \item $p^* = 8$ y $T^* = 8$
    \item $p^* = 5$ y $T^* = 24.5$
\end{opts}

\q{Un monopolista produce dos bienes digitales con costo marginal nulo para dos consumidores cuyas valoraciones $(v_A, v_B)$ son $(100, 40)$ y $(40, 100)$. ¿Cuáles son los ingresos máximos con venta separada y con empaquetamiento puro, respectivamente?}
\begin{opts}
    \item $160$ y $280$
    \item $200$ y $280$
    \item $200$ y $240$
    \item $180$ y $260$
\end{opts}

\q{Un monopolista con $CMg_A = CMg_B = 20$ enfrenta tres consumidores con valoraciones $(v_A, v_B)$ iguales a $(90, 10)$, $(60, 60)$ y $(10, 90)$. Si aplica empaquetamiento mixto con $P_A = 90$, $P_B = 90$ y $P_{AB} = 120$, ¿cuál es su beneficio total?}
\begin{opts}
    \item $180$ unidades monetarias
    \item $160$ unidades monetarias
    \item $220$ unidades monetarias
    \item $200$ unidades monetarias
\end{opts}

\q{Una empresa de servicios públicos enfrenta demanda pico $P_1 = 100 - Q_1$ y demanda valle $P_2 = 40 - Q_2$. El costo marginal operativo es $b = 10$ y el costo marginal de capacidad es $\beta = 20$. ¿Cuáles son los precios eficientes de Steiner-Boiteux y la capacidad óptima $K^*$?}
\begin{opts}
    \item $P_1^* = 20$, $P_2^* = 20$, $K^* = 80$
    \item $P_1^* = 30$, $P_2^* = 20$, $K^* = 70$
    \item $P_1^* = 25$, $P_2^* = 15$, $K^* = 75$
    \item $P_1^* = 30$, $P_2^* = 10$, $K^* = 70$
\end{opts}

\q{Un monopolio natural tiene función de costos $CT(Q) = 400 + 10Q$ y enfrenta la demanda $P = 90 - Q$. Si el regulador impone eficiencia asignativa ($P = CMg$), ¿cuál es la cantidad producida y la transferencia fiscal necesaria para evitar pérdidas?}
\begin{opts}
    \item $Q^* = 80$ y transferencia de $400$
    \item $Q^* = 40$ y transferencia de $400$
    \item $Q^* = 80$ y transferencia de $800$
    \item $Q^* = 70$ y transferencia de $300$
\end{opts}

\q{Un monopolio natural presenta costos $CT(Q) = 400 + 10Q$ y demanda $P = 60 - Q$. Si el regulador fija el precio de segundo mejor igual al costo medio ($P = CMe$) maximizando el excedente del consumidor, ¿cuáles son $Q^*$ y $P^*$?}
\begin{opts}
    \item $Q^* = 10$ y $P^* = 50$
    \item $Q^* = 40$ y $P^* = 20$
    \item $Q^* = 25$ y $P^* = 35$
    \item $Q^* = 50$ y $P^* = 10$
\end{opts}

\q{Una industria tiene función de costos $C(Q) = 50 + 2Q^2$ (cuyo costo medio mínimo ocurre en $Q = 5$). Comparando una sola empresa frente a dos empresas simétricas, ¿cuál afirmación es correcta para $Q = 6$ y $Q = 8$?}
\begin{opts}
    \item No es subaditiva en $Q = 6$ ni en $Q = 8$ porque ambos superan el mínimo del costo medio.
    \item Es subaditiva tanto en $Q = 6$ como en $Q = 8$ debido a la presencia del costo fijo de $50$.
    \item Es subaditiva en $Q = 6$ ($122 < 136$), pero deja de serlo en $Q = 8$ ($178 > 164$).
    \item Es subaditiva en $Q = 8$ ($164 < 178$), pero no en $Q = 6$ ($136 > 122$).
\end{opts}

\q{Una empresa regulada provee dos servicios independientes con elasticidades $\varepsilon_1 = -1.5$ y $\varepsilon_2 = -3.0$, y costos marginales $CMg_1 = 12$ y $CMg_2 = 9$. Si el número de Ramsey es $R = \frac{\lambda}{1+\lambda} = 0.30$, ¿cuáles son los precios de Ramsey $(P_1^*, P_2^*)$?}
\begin{opts}
    \item $P_1^* = 14.40$ y $P_2^* = 9.90$
    \item $P_1^* = 16.00$ y $P_2^* = 12.00$
    \item $P_1^* = 13.50$ y $P_2^* = 11.25$
    \item $P_1^* = 15.00$ y $P_2^* = 10.00$
\end{opts}

\q{Una EPS regulada bajo un esquema de precio tope $P_t = P_{t-1}(1 + \pi - X)$ parte de una tarifa media de $S/\ 4.00/\text{m}^3$. Si la inflación anual es $\pi = 5\%$ y el factor de productividad fijado por el regulador es $X = 2\%$, ¿cuál es la tarifa máxima permitida para el periodo $t$?}
\begin{opts}
    \item $S/\ 4.12/\text{m}^3$
    \item $S/\ 4.20/\text{m}^3$
    \item $S/\ 4.28/\text{m}^3$
    \item $S/\ 3.92/\text{m}^3$
\end{opts}

\q{Una empresa regulada por costo de servicio tiene una base de capital $K = 2000$, una tasa de retorno permitida $s = 8\%$, gastos operativos de $400$ y depreciación anual de $100$. Si el volumen facturado es $Q = 220$, ¿cuál es la tarifa media requerida?}
\begin{opts}
    \item $2.27$ por unidad
    \item $3.00$ por unidad
    \item $2.55$ por unidad
    \item $3.50$ por unidad
\end{opts}

\q{Una empresa multiproducto presenta costos stand-alone $C(q_1, 0) = 100$ y $C(0, q_2) = 80$, y un costo conjunto $C(q_1, q_2) = 150$. Si cubre exactamente sus costos totales, ¿en qué intervalo debe ubicarse el ingreso $R_1$ del servicio 1 para estar libre de subsidios cruzados?}
\begin{opts}
    \item $50 \le R_1 \le 100$
    \item $70 \le R_1 \le 150$
    \item $70 \le R_1 \le 100$
    \item $80 \le R_1 \le 100$
\end{opts}

\q{La estructura tarifaria de agua potable de una EPS contempla un cargo fijo de $S/\ 5.00$, un primer bloque de $0$ a $10\ \text{m}^3$ a $S/\ 1.50/\text{m}^3$ y un segundo bloque de más de $10$ hasta $25\ \text{m}^3$ a $S/\ 3.00/\text{m}^3$. ¿Cuánto pagará un hogar que consume $20\ \text{m}^3$?}
\begin{opts}
    \item $S/\ 65.00$
    \item $S/\ 45.00$
    \item $S/\ 60.00$
    \item $S/\ 50.00$
\end{opts}

\q{Un monopolista con $CMg = 20$ enfrenta la demanda $P = 100 - Q$. Si el regulador fija un precio máximo $P_{\max} = 40$, ¿cuál es la cantidad ofrecida $Q_{\text{reg}}$ y cuánto se reduce la pérdida irrecuperable de eficiencia?}
\begin{opts}
    \item $Q_{\text{reg}} = 60$ y la PBE se reduce en $600$
    \item $Q_{\text{reg}} = 60$ y la PBE se reduce en $200$
    \item $Q_{\text{reg}} = 40$ y la PBE se reduce en $400$
    \item $Q_{\text{reg}} = 80$ y la PBE se reduce en $800$
\end{opts}

\q{Un monopolista elige simultáneamente precio y gasto en publicidad $A$. Si en el óptimo la elasticidad precio de la demanda es $\varepsilon_P = -2.5$ y la elasticidad de la demanda respecto a la publicidad es $\varepsilon_A = 0.5$, ¿cuál es la razón óptima gasto en publicidad sobre ingresos $A / (PQ)$ según Dorfman-Steiner?}
\begin{opts}
    \item $0.10$ (equivalente al $10\%$)
    \item $0.20$ (equivalente al $20\%$)
    \item $0.25$ (equivalente al $25\%$)
    \item $0.50$ (equivalente al $50\%$)
\end{opts}
